\section{Proactive Memory}
\label{sec:framework}

\system{} uses three complementary memory layers (Figure~\ref{fig:architecture}): event memory for recent public events, temporal trajectory memory for source-checked activity summaries, and behavioral regularity memory for replay-assessed conditional regularities.
Construction checks source fidelity and historical support.
The decision stage prepares memory and matches current conditions before semantic screening; contextual assessment combines current and historical evidence.

\begin{figure*}[t]
    \centering
    \makebox[\textwidth][c]{\resizebox{0.98\textwidth}{!}{% Editable flat vector architecture, matching the visual language of Figure 1.
% Rows align each memory layer with the module that builds it (left) and the
% place that reads it (right): S_t and W_t enter the context directly, while
% L_t is matched against the current event by retrieval. Example values follow
% the Figure 1 walking scenario and are illustrative.
\definecolor{archInk}{HTML}{2B3A4A}
\definecolor{archMuted}{HTML}{6E7C8A}
\definecolor{archLaneA}{HTML}{F2F5F8}
\definecolor{archLaneM}{HTML}{EAF5F1}
\definecolor{archLaneB}{HTML}{FDF5EA}
\definecolor{archBlue}{HTML}{5B7FA3}
\definecolor{archBlueLt}{HTML}{DCE6EF}
\definecolor{archTeal}{HTML}{2A8C7E}
\definecolor{archTealLt}{HTML}{D3EBE4}
\definecolor{archAmber}{HTML}{E0913A}
\definecolor{archBrown}{HTML}{A8692A}
\definecolor{archPurple}{HTML}{7B6BB5}
\definecolor{archPurpleLt}{HTML}{E6E1F4}
\definecolor{archRose}{HTML}{C75B66}
\definecolor{archGreyLt}{HTML}{E4E9EE}
\definecolor{archPanel}{HTML}{F7F9FB}
\begin{tikzpicture}[x=1cm,y=-1cm,line cap=round,line join=round,
  font=\sffamily\fontsize{8.5}{10}\selectfont,text=archInk,
  txt/.style={inner sep=0pt,align=left,text=archInk},
  sm/.style={txt,font=\sffamily\fontsize{7.2}{8.6}\selectfont},
  ex/.style={txt,font=\sffamily\fontsize{6.6}{7.9}\selectfont},
  hd/.style={txt,font=\sffamily\bfseries\fontsize{9.5}{11}\selectfont},
  ttl/.style={txt,font=\sffamily\bfseries\fontsize{8.5}{10}\selectfont},
  card/.style={fill=white,rounded corners=4pt},
  inset/.style={fill=archPanel,rounded corners=3pt},
  arr/.style={-{Stealth[length=1.9mm,width=1.4mm]},line width=1pt,rounded corners=3pt},
  lab/.style={sm,fill=white,inner sep=1.2pt,rounded corners=1pt}]
\path[use as bounding box] (0,0) rectangle (16.80,9.30);

% ---------------------------------------------------------------- icons
\newcommand{\llmicon}[3]{% x, y, colour
  \fill[#3,rounded corners=2.5pt] (#1-.26,#2-.20) rectangle (#1+.26,#2+.20);
  \node[txt,font=\sffamily\bfseries\fontsize{6}{7}\selectfont,text=white] at (#1,#2) {LLM};}
\newcommand{\eventsicon}[2]{% x, y
  \foreach \dx/\dy in {-.13/-.13,.13/-.13,-.13/.13}{
    \fill[archBlue,rounded corners=.8pt] (#1+\dx-.10,#2+\dy-.10) rectangle ++(.20,.20);}
  \fill[archAmber,rounded corners=.8pt] (#1+.03,#2+.03) rectangle ++(.20,.20);}
\newcommand{\agenticon}[3]{% x, y, colour
  \begin{scope}[shift={(#1,#2)},scale=.95]
    \draw[#3,line width=1pt] (0,-.25)--(0,-.34);
    \fill[#3] (0,-.37) circle (.045);
    \fill[#3,rounded corners=1pt] (-.35,-.12) rectangle (-.28,.04);
    \fill[#3,rounded corners=1pt] (.28,-.12) rectangle (.35,.04);
    \fill[#3,rounded corners=3.5pt] (-.29,-.25) rectangle (.29,.17);
    \fill[white,rounded corners=2.5pt] (-.20,-.15) rectangle (.20,.08);
    \fill[#3] (-.085,-.055) circle (.042) (.085,-.055) circle (.042);
    \draw[#3,line width=.8pt] (-.055,.01) .. controls (-.02,.045) and (.02,.045) .. (.055,.01);
  \end{scope}}
\newcommand{\filtericon}[3]{% funnel: x, y, colour
  \fill[#3] (#1-.24,#2-.20)--(#1+.24,#2-.20)--(#1+.06,#2+.02)--(#1+.06,#2+.20)
    --(#1-.06,#2+.13)--(#1-.06,#2+.02)--cycle;}
\newcommand{\magicon}[3]{% x, y, colour
  \draw[#3,line width=1.4pt,fill=white] (#1-.04,#2-.04) circle (.13);
  \draw[#3,line width=2pt] (#1+.06,#2+.06)--(#1+.18,#2+.18);}
\newcommand{\ok}[3]{\draw[#3,line width=.95pt] (#1-.06,#2)--(#1-.015,#2+.05)--(#1+.07,#2-.06);}
\newcommand{\no}[3]{\draw[#3,line width=.95pt] (#1-.05,#2-.05)--(#1+.05,#2+.05) (#1-.05,#2+.05)--(#1+.05,#2-.05);}
% memory-cylinder layer: y-top, y-bottom, colour
\newcommand{\cyl}[3]{%
  \fill[#3] (6.20,#1) rectangle (8.30,#2);
  \fill[#3] (7.25,#2) ellipse (1.05 and .18);
  \draw[white,line width=.9pt] (6.20,#1) arc[start angle=180,end angle=360,x radius=1.05,y radius=.18];}
% replay day: x, y, letter, state (1 fulfilled, 2 violated, 3 excluded)
\newcommand{\rday}[4]{%
  \ifcase#4\or
    \fill[archTeal] (#1,#2) circle (.15);
    \node[txt,font=\sffamily\bfseries\fontsize{6}{7}\selectfont,text=white] at (#1,#2) {#3};\or
    \fill[archRose] (#1,#2) circle (.15);
    \node[txt,font=\sffamily\bfseries\fontsize{6}{7}\selectfont,text=white] at (#1,#2) {#3};\or
    \draw[archMuted!60,line width=.7pt,fill=white] (#1,#2) circle (.14);
    \node[txt,font=\sffamily\fontsize{6}{7}\selectfont,text=archMuted] at (#1,#2) {#3};\fi}

% ---------------------------------------------------------------- lanes
\fill[archLaneA,rounded corners=6pt] (.15,.10) rectangle (5.85,8.65);
\fill[archLaneM,rounded corners=6pt] (6.00,.10) rectangle (8.50,8.65);
\fill[archLaneB,rounded corners=6pt] (8.65,.10) rectangle (16.65,8.65);
\node[hd,text=archBlue,anchor=west] at (.35,.40) {A\enspace Memory construction};
\node[hd,text=archTeal,anchor=center] at (7.25,.40) {ProactMem};
\node[hd,text=archBrown,anchor=west] at (8.85,.40) {B\enspace Event-driven decision};
\node[sm,text=archMuted,anchor=east] at (16.50,.40) {online, per event};

% ================================================================ A
% ---- row 1: committed events
\fill[card] (.35,.85) rectangle (5.65,2.85);
\eventsicon{.78}{1.18}
\node[ttl,anchor=west] at (1.20,1.12) {Committed events};
\node[sm,text=archMuted,anchor=west] at (1.20,1.47) {daily buffer of the open local day};
\fill[inset] (.55,1.78) rectangle (5.45,2.68);
\draw[archMuted!50,line width=.8pt] (.80,2.08)--(5.20,2.08);
\foreach \x/\t/\c in {.95/17:41/archBlue!45,1.75/17:55/archBlue!45,2.55/18:22/archTeal,
                       3.35/18:48/archTeal,4.15/19:30/archBlue!45}{
  \fill[\c,rounded corners=.8pt] (\x-.10,1.98) rectangle (\x+.10,2.18);
  \node[ex,text=archMuted,anchor=center] at (\x,2.42) {\t};}
\node[ex,text=archMuted,anchor=center] at (4.85,2.08) {$\cdots$};
% ---- row 2: trajectory extraction
\fill[card] (.35,3.35) rectangle (5.65,5.55);
\llmicon{.78}{3.68}{archTeal}
\node[ttl,anchor=west] at (1.20,3.62) {Trajectory extraction};
\node[sm,text=archMuted,anchor=west] at (1.20,3.97) {segments with typed facts and citations};
\fill[inset] (.55,4.25) rectangle (5.45,5.38);
\node[ex,anchor=west] at (.70,4.43) {\textbf{18:22--18:48 walking}\enspace\textcolor{archMuted}{steps 2{,}310\enspace[$e_{17},e_{18}$]}};
\ok{.78}{4.76}{archTeal}
\node[ex,text=archTeal,anchor=west] at (.92,4.76) {facts $\subseteq$ cited events $\Rightarrow$ kept in $W_t$};
\no{.78}{5.09}{archRose}
\node[ex,text=archRose,anchor=west] at (.92,5.09) {``walked 5\,km'' has no source fact $\Rightarrow$ rejected};
% ---- row 3: regularity proposal and replay
\fill[card] (.35,6.05) rectangle (5.65,8.45);
\llmicon{.78}{6.38}{archPurple}
\node[ttl,anchor=west] at (1.20,6.32) {Regularity proposal};
\node[sm,text=archMuted,anchor=west] at (1.20,6.67) {weekly, over 7 completed records};
\fill[inset] (.55,6.93) rectangle (5.45,8.30);
\node[ex,anchor=west] at (.70,7.14) {\textcolor{archPurple}{\textbf{trigger}} weekday, low activity after 17:30};
\node[ex,anchor=west] at (.70,7.42) {\textcolor{archPurple}{\textbf{expect}}\ \ walk within 18:00--19:00};
\node[ex,text=archMuted,anchor=west] at (.70,7.78) {replay};
\foreach \x/\d/\s in {1.62/M/1,1.98/T/1,2.34/W/2,2.70/T/1,3.06/F/1,3.42/S/3,3.78/S/3}{\rday{\x}{7.78}{\d}{\s}}
\node[ex,text=archMuted,anchor=west] at (4.02,7.78) {weekend excl.};
\node[ex,anchor=west] at (.70,8.10) {$n_f{=}4,\ n_v{=}1$, conf $0.80$, lift $\ge1$\enspace$\Rightarrow$\enspace\textcolor{archPurple}{\textbf{active}}};
% vertical flow
\draw[arr,draw=archMuted] (3.00,2.85)--(3.00,3.35);
\node[sm,text=archMuted,anchor=west] at (3.12,3.10) {day closes};
\draw[arr,draw=archMuted] (3.00,5.55)--(3.00,6.05);
\node[sm,text=archMuted,anchor=west] at (3.12,5.80) {7 daily records};
% writes into the memory layers
\draw[arr,draw=archBlue] (5.65,1.85)--(6.15,1.85);
\draw[arr,draw=archTeal] (5.65,4.45)--(6.15,4.45);
\draw[arr,draw=archPurple] (5.65,7.15)--(6.15,7.15);

% ================================================================ memory
\cyl{5.80}{8.25}{archPurple}
\cyl{3.10}{5.80}{archTeal}
\cyl{.95}{3.10}{archBlue}
\fill[archBlue!55] (7.25,.95) ellipse (1.05 and .18);
\newcommand{\cyltext}[4]{% y, title, symbol, note
  \node[txt,font=\sffamily\bfseries\fontsize{8}{9.6}\selectfont,text=white,align=center] at (7.25,#1) {#2};
  \node[txt,font=\fontsize{9}{10}\selectfont,text=white] at (7.25,#1+.42) {#3};
  \node[txt,font=\sffamily\fontsize{6.6}{8}\selectfont,text=white,align=center] at (7.25,#1+.82) {#4};}
\cyltext{1.55}{Event\\memory}{$S_t$}{last 72\,h,\\[-1pt]$\le$99 events}
\cyltext{4.05}{Temporal\\trajectories}{$W_t$}{source-checked\\[-1pt]segments}
\cyltext{6.80}{Behavioral\\regularities}{$L_t$}{active: visible\\[-1pt]candidate: hidden}

% ================================================================ B
% ---- retrieval of applicable regularities (row 3)
\fill[card] (8.85,6.05) rectangle (11.55,7.85);
\magicon{9.20}{6.40}{archAmber}
\node[ttl,anchor=west] at (9.50,6.32) {Retrieval};
\node[sm,text=archMuted,anchor=east] at (11.42,6.34) {top 3};
\foreach \y/\t in {6.78/{active regularity},7.12/{trigger + time match},7.46/{not yet fulfilled}}{
  \ok{9.10}{\y}{archTeal}
  \node[ex,anchor=west] at (9.25,\y) {\t};}
% ---- current event
\fill[archAmber,rounded corners=4pt] (8.85,8.02) rectangle (11.55,8.45);
\node[txt,font=\sffamily\bfseries\fontsize{6.8}{8}\selectfont,text=white] at (10.20,8.235) {$e_t$\ \ 18:30, low activity};
\draw[arr,draw=archAmber] (10.20,8.02)--(10.20,7.85);
% ---- memory reads
\draw[arr,draw=archBlue] (8.32,1.85)--(11.90,1.85);
\node[lab,text=archBlue] at (10.10,1.85) {$S_t$};
\draw[arr,draw=archTeal] (8.32,4.45)--(11.90,4.45);
\node[lab,text=archTeal] at (10.10,4.45) {$W_t$};
\draw[arr,draw=archPurple] (8.32,6.95)--(8.85,6.95);
\draw[arr,draw=archPurple] (11.55,6.95)--(11.90,6.95);
\node[sm,text=archPurple,anchor=south] at (11.72,6.82) {$R_t$};

% ---- context C_t
\fill[card] (11.90,.85) rectangle (13.90,8.45);
\node[ttl,anchor=center] at (12.90,1.15) {Context $C_t$};
\newcommand{\ctxchip}[4]{% y, fill, text colour, label
  \fill[#2,rounded corners=2.5pt] (12.05,#1-.21) rectangle (13.75,#1+.21);
  \node[sm,text=#3,anchor=center] at (12.90,#1) {#4};}
\ctxchip{1.85}{archBlueLt}{archBlue}{$S_t$\, recent events}
\ctxchip{3.10}{archGreyLt}{archMuted}{$Q_t$\, episode state}
\ctxchip{4.45}{archTealLt}{archTeal}{$W_t$\, trajectories}
\ctxchip{6.95}{archPurpleLt}{archPurple}{$R_t$\, regularities}
\ctxchip{8.10}{archAmber!25}{archBrown}{$e_t$\, current event}
\draw[arr,draw=archAmber] (11.55,8.235)--(11.90,8.235);
\node[ex,text=archMuted,align=center,anchor=center] at (12.90,5.70) {$C_t=(e_t,S_t,$\\$W_t,R_t,Q_t)$};

% ---- decision LLM (row 1)
\draw[arr,draw=archInk] (13.90,1.85)--(14.25,1.85);
\fill[card] (14.25,.85) rectangle (16.45,2.85);
\agenticon{14.68}{1.22}{archTeal}
\node[ttl,anchor=west] at (15.08,1.18) {Decision};
\node[ttl,anchor=west] at (15.08,1.50) {LLM};
\node[ex,anchor=west] at (14.42,1.98) {$(\tilde y_t,\tilde m_t)=\pi_\theta(C_t)$};
\node[ex,text=archMuted,anchor=west] at (14.42,2.30) {timely, feasible,};
\node[ex,text=archMuted,anchor=west] at (14.42,2.56) {low-risk service?};
% ---- intervention control (row 2)
\draw[arr,draw=archInk] (15.35,2.85)--(15.35,3.35);
\fill[card] (14.25,3.35) rectangle (16.45,5.55);
\filtericon{14.65}{3.72}{archBrown}
\node[ttl,anchor=west] at (15.05,3.62) {Intervention};
\node[ttl,anchor=west] at (15.05,3.94) {control};
\node[ex,text=archMuted,anchor=west] at (14.42,4.42) {compare with allowed};
\node[ex,text=archMuted,anchor=west] at (14.42,4.68) {interventions;};
\no{14.50}{5.05}{archRose}
\node[ex,anchor=west] at (14.62,5.05) {duplicate $\to$ abstain};
% ---- output (row 3)
\draw[arr,draw=archInk] (15.35,5.55)--(15.35,6.05);
\fill[archTeal,rounded corners=4pt] (14.25,6.05) rectangle (16.45,7.55);
\node[txt,font=\sffamily\bfseries\fontsize{7.4}{8.8}\selectfont,text=white,anchor=center] at (15.35,6.34) {\textsc{Intervene}};
\node[txt,font=\sffamily\itshape\fontsize{7}{8.4}\selectfont,text=white,align=center,anchor=center]
  at (15.35,6.98) {``Up for a short\\walk?''};
\fill[archGreyLt,rounded corners=4pt] (14.25,7.68) rectangle (16.45,8.45);
\node[txt,font=\sffamily\bfseries\fontsize{7.4}{8.8}\selectfont,text=archMuted,anchor=center] at (15.35,7.98) {\textsc{Abstain}};
\node[ex,text=archMuted,anchor=center] at (15.35,8.25) {otherwise};

% ---------------------------------------------------------------- feedback
\draw[arr,draw=archTeal,dash pattern=on 2.4pt off 1.8pt]
  (16.45,7.615)--(16.72,7.615)--(16.72,8.95)--(.05,8.95)--(.05,1.85)--(.35,1.85);
\node[lab,text=archTeal,fill=white] at (7.25,8.95) {commit $e_t$ after the decision is finalized; it affects only later decisions};
\end{tikzpicture}
}}
    \caption{\system{} architecture. (A) Source validation and replay build hierarchical memory. (B) Memory preparation and matching precede screening, which routes events to abstention or contextual assessment with intervention control. Both routes commit the event after finalization (dashed arrows).}
    \label{fig:architecture}
\end{figure*}

\subsection{Memory construction and update}

Service opportunities can depend on both recent changes and recurring behavior, motivating a memory hierarchy that preserves fine-grained evidence while validating longer-term abstractions.

\paragraph{bounded event memory.}
Evidence utilization can degrade in long contexts~\citep{liu2024lost}, motivating a bounded event layer that preserves recent observations at their original granularity.
For event $e_t$, event memory retains prior events in their original public representation, intersecting a temporal horizon $\Delta$ with capacity $K$:
\begin{equation}
S_t=\operatorname{Tail}_{K-1}
\left((e_i)_{i<t,\;0\leq\tau_t-\tau_i\leq\Delta}\right).
\label{eq:source_memory}
\end{equation}
Here $\operatorname{Tail}_{K-1}$ selects the most recent $K-1$ events in the filtered sequence.
We use $\Delta=72$ hours and $K=100$: at most 99 prior events accompany the separately supplied current event.
A separate buffer retains all committed events from the open local day for trajectory construction.

\paragraph{fact-validated temporal trajectory memory.}
To retain activity context beyond the recent window while keeping abstractions traceable, we construct temporal trajectories with source-level fact validation.
After a local day closes, a construction LLM summarizes its buffered events as trajectory segments with intervals, typed facts, and event citations.
Let $F(\cdot)$ denote normalized checkable facts obtained from typed fields.
For segment $d$ and its cited events $\mathcal{E}(d)$, source consistency requires
\begin{equation}
F(d)\subseteq\bigcup_{e\in\mathcal{E}(d)}F(e),
\qquad \mathcal{E}(d)\subseteq E^u_{\text{completed day}}.
\label{eq:daily_validation}
\end{equation}
Here $E^u_{\text{completed day}}$ contains the same entity's events from that completed day.
The validator checks entity, interval, citations, and canonicalized scalar values at compatible field paths.
Acceptance requires source-supported typed facts and resolved citations.
Accepted temporal trajectories form $W_t$, retaining historical context beyond event memory.
Appendix~\ref{app:implementation} details field handling and buffering.

\paragraph{replay-assessed behavioral regularity memory.}
To ground behavioral expectations in observed recurrence, we assess proposed conditional regularities through historical replay and revise their eligibility as evidence accumulates.
Weekly consolidation groups seven completed daily records per entity in time order, with each record assigned to one cycle.
A cycle may span more than seven calendar days.
An LLM proposes a trigger $g$, expected behavior $p$, time window, optional weekday restrictions, and source daily identifiers.
Replay checks these conditions against validated daily facts: an observable triggered day is fulfilled when $p$ occurs at or after the trigger within the window and violated otherwise.
Missing observations are censored, inapplicable days are excluded, and each date contributes at most once.
Let $n_f$ count fulfilled (supporting) days and $n_v$ count violated days:
\begin{equation}
\operatorname{conf}_{\mathrm{replay}}=\frac{n_f}{n_f+n_v},
\qquad
\operatorname{lift}=\frac{\operatorname{conf}_{\mathrm{replay}}}{\operatorname{base}(p)}.
\label{eq:replay_gate}
\end{equation}
The base rate is the frequency of $p$ over observable days under the same time-window and weekday restrictions across both triggered and other observable days.
A zero denominator gives a zero score.

A well-formed proposal becomes \emph{active} when $n_f\geq3$, confidence $\geq0.5$, and lift $\geq1.0$; otherwise it remains a decision-invisible \emph{candidate}.
Malformed proposals are rejected.
Consolidation resets the online trajectory view, retaining records for audit.
Later observations update regularity confidence and lifecycle: repeated violations can lower retrieval priority or suspend eligibility.
Appendix~\ref{app:lifecycle} specifies these transitions.

\subsection{Memory retrieval and selection}

To make accumulated regularities relevant to the present service opportunity, retrieval checks their trigger conditions, temporal applicability, and occurrence state.
The event memory supplies the bounded sequence and $W_t$ supplies all completed temporal trajectories in the current cycle.
Retrieval of regularities additionally checks entity-local trigger context, weekdays, and temporal applicability with a 30-minute lead and grace interval.
Matching considers current and previously recorded causal trigger occurrences.
An occurrence is fulfilled when its expected behavior has been observed; censored occurrences have insufficient observations for assessment.
Selection excludes both states.
For regularity store $L_t^u$, lifecycle eligibility $\operatorname{Eligible}_t$, trigger and time matching $\operatorname{Match}_t$, and occurrence state $s_t$,
\begin{equation}
\begin{aligned}
R_t&=\operatorname{Top3}\big\{\ell\in L_t^u:\\
&\quad\operatorname{Eligible}_t(\ell)\land\operatorname{Match}_t(\ell),\\
&\quad s_t(\ell)\notin\{\text{fulfilled},\text{censored}\}\big\}.
\end{aligned}
\label{eq:long_term_retrieval}
\end{equation}
The top-three selection ranks eligible regularities by occurrence state, lifecycle, confidence, and temporal distance (Appendix~\ref{app:matching}).
Retrieval precedes screening and supplies applicable regularities and occurrence states.

\subsection{Proactive decision}
\label{sec:decision_control}

Timely assistance requires combining current needs with historical context~\citep{yang2025contextagent}.
The decision stage screens observations, assesses service value, and tracks prior opportunity coverage.

\paragraph{Semantic screening.}
\label{sec:semantic_screening}
Screening uses prepared memory and current-condition matches to assess relevance and event-local abstention evidence.
Let $\mathcal{M}_t^{(k)}$ contain prior same-entity event records, validated trajectory segments, and active or weakening regularities for $k=1,2,3$, respectively.
For a nonempty layer, semantic encoder $\phi$ gives the maximum cosine similarity:
\begin{equation}
\rho_t^{(k)}=\max_{m\in\mathcal{M}_t^{(k)}}
\cos\!\left(\phi(e_t),\phi(m)\right).
\label{eq:semantic_screening}
\end{equation}
Each layer has a relevance threshold $\lambda_k$.
Pending services, timer observations, and matched upcoming or unresolved regularity occurrences take the \emph{escalate} route, as do incomplete memory and any $\rho_t^{(k)}\geq\lambda_k$.
When all three populated layers satisfy $\rho_t^{(k)}<\lambda_k$, a single-event LLM assesses the public current event.
A valid \emph{skip} requires explicit abstention evidence quoted from the event and reported confidence at least $\eta$.
All other outcomes, including new needs, risks, ambiguity, and history-dependent judgments, escalate.
An accepted skip yields \emph{abstain}; escalated events receive contextual assessment.
Appendix~\ref{app:semantic_screening} specifies the screening contract and execution settings.

\paragraph{Contextual assessment and intervention control.}
For escalated events, the decision context combines the current event, recent events, temporal trajectories, selected regularities, and observable episode state:
\begin{equation}
C_t=(e_t,S_t,W_t,R_t,Q_t),
\label{eq:decision_context}
\end{equation}
where $Q_t$ describes the service opportunity's observable state; the controller separately records previously allowed interventions to track service coverage.
The LLM produces $(\tilde y_t,\tilde m_t)=\pi_\theta(C_t)$, proposing intervention when the current event \emph{or} historical context supports concrete, timely, feasible, low-risk assistance, and abstention otherwise.

Within each scope, all methods apply the same intervention controller policy to their own context and intervention history.
The controller compares proposed services with previously allowed interventions and converts duplicates to \emph{abstain}.
Scope-specific changes in opportunity state, service identity, contextual regularities, or local day permit renewed assessment; each intervention originates from an LLM proposal.
The controller finalizes the proposal as $(\hat y_t,\hat m_t)$, the output defined in Section~\ref{sec:problem}.
Appendix~\ref{app:intervention_control} specifies matching fields, phase transitions, and re-entry rules.

Inference stages $e_t$, retrieves prior memory and matches current conditions, screens the event, assembles full context for escalated events, finalizes the output, and commits $e_t$.
Both routes retain the complete public event in the event buffer and open-day buffer, advancing trajectory construction, regularity replay, and lifecycle updates.
Maintenance incorporating $e_t$ can affect only later decisions.
Reference labels remain private during construction and retrieval.
